According to the press reports consulted, on 8 October 2026 a mudflow descended from the San Carlos de Apoquindo ravine (Las Condes, Santiago) over a hillside burned nine months earlier \citep{biobio2026origen}. The flow was modeled using a two-dimensional shallow-water model with mud rheology, on terrain with a resolution of \SI{30}{m} resampled to a \SI{4}{m} grid with 42,018 buildings, using only public data and an evidence log of every source consulted. A total of 84 simulations were completed: 19 screening runs, 41 full-city runs, 5 quadrature runs and 19 campus micro-simulations. With about \SI{144}{mm} of accumulated rainfall as reported by the mayor \citep{interferencia2026}, and with the \SI{0}{\celsius} isotherm reported by CR2 above the summit of the catchment \citep{cr2_2026} (calculated area: \SI{10.1}{km^2}), an independent estimate adopting the reported soil saturation places the mixture volume at about \SI{295000}{m^3}, of the same order of magnitude as that required by the simulated reach. The hypothesis selected for the held-out test (z02)---main channel, southern branches and rainfall on the hillside---reproduces the Camino El Alba, General Blanche and Quebrada Honda corridors that delimit the official affected area and the reach to Padre Hurtado, and calculates about \SI{1}{m} of mud next to Clínica UC Christus, with an exploratory compatibility index (ECI) of 72/100. This index compares hypotheses and does not measure predictive accuracy. A test using held-out observations, defined before examining results, had an unfavorable outcome: the simulated flow reaches 19 of 49 officially affected facilities (\SI{39}{\percent}, compared with the specified \SI{70}{\percent}), because the model leaves parts of the official affected area without mud. In subsequent fitting tests, mud entry into the represented houses and the assumption of turbid water on campus and on La Plaza Sur, together with mud in the main branch, raise the ECI verified with a fine time step to 76--82/100. The integrated hypothesis obtains 79--82/100: 82 with assumed rainfall of \SI{35}{mm/h} and a fine time step (z32), 81 with the recorded Río Molina rainfall pulse (p34b, corrected) and 79 with that rainfall distributed over \SI{4}{h} (v33). Burn severity estimated with Sentinel-2 shows a burn signal in \SI{82}{\percent} of the catchments, mostly at low or moderate-low severity. With the coefficients assumed for the scenario without fire, the balance calculates about \SI{62}{\percent} of the mud volume and \SI{43}{\percent} of the solids of the scenario with fire. In the model, an assumed log obstacle at the end of the main channel does not divert the flow southward. In hypothetical scenarios, flooding of Quebrada Honda is reproduced if the sink representing the Francisco Bulnes Correa culvert does not remove all the flow calculated at its inlet, of about \SI{65}{m^3/s}. An assumed capacity of \SI{30}{m^3/s} leaves surface flow. These values do not describe the real structure. Retention in the represented houses improves agreement with the official affected area, but no configuration reproduces both that extent and the meter reported at Clínica UC Christus \citep{cadena3_2026}. The micro-simulations and the city simulation with a fine mesh on campus (m36) reproduce entry through the clearing into the sunken courtyard; m36 calculates \SI{0.17}{m} in the courtyard and wets \SI{58}{\percent} of La Plaza Sur. The study closed on 13 October 2026; its fitting results have no independent validation.
